\documentclass[11pt]{article}
\usepackage[margin=1in]{geometry}
\usepackage[T1]{fontenc}
\usepackage[utf8]{inputenc}
\usepackage{amsmath,amssymb,amsthm}
\usepackage{mathtools}
\usepackage{listings}
\usepackage{hyperref}

\lstset{
  basicstyle=\ttfamily\small,
  columns=fullflexible,
  keepspaces=true,
  showstringspaces=false,
  frame=single
}

\newtheorem{definition}{Definition}
\newtheorem{proposition}{Proposition}

\newcommand{\Types}{\mathcal{T}}
\newcommand{\Effects}{\mathcal{E}}
\newcommand{\glb}{\mathbin{\sqcap}}
\newcommand{\compose}{\mathbin{\cdot}}
\newcommand{\eps}{\varepsilon}
\newcommand{\id}{\mathbf{1}}
\newcommand{\sem}[1]{[\![#1]\!]}

\title{A Declarative Stack-Effect Type Checker for a Forth-like Language}
\author{Jaanus Pöial \\ Tallinn University of Technology}
\date{}

\begin{document}
\maketitle

\begin{abstract}
This paper describes a static type checker for a Forth(-like) language whose semantics are expressed as typed stack effects. The checker is implemented as a single Gforth program, but its core ideas are language independent: a finite type lattice is loaded from an external file, word specifications are given declaratively, and user-defined words are inferred by composing abstract stack transformations instead of executing the program. The central technical device is a partial algebra of stack effects. Sequential composition unifies adjacent stack symbols by taking the greatest lower bound of comparable types; branching computes a greatest lower bound of whole effects; and iteration is accepted only when repeated execution stabilizes to an idempotent effect. The result is a compact checker that supports ordinary words, literals, parser words, defining words, and declaratively specified control structures such as \texttt{IF ... ELSE ... FI} and \texttt{BEGIN ... WHILE ... REPEAT}. We present the mathematical model, the operations implemented by the checker, and worked examples taken from the bundled specification and program files.
\end{abstract}

\section{Introduction}
Concatenative languages place computation on the operand stack. A word is naturally described by the stack cells it consumes and the cells it produces, so the obvious static abstraction is a \emph{stack effect}. In simple cases, informal comments such as \texttt{( n n -- n )} are already useful documentation in Forth practice and standardization documents \cite{rather1996evolution,forth2012standard}. The goal of the checker described here is to turn such annotations into an executable static semantics.

The checker takes three external inputs. A \texttt{types} file defines a finite subtype hierarchy. A \texttt{specs} file assigns effects and parsing behavior to words, literals, defining words, and control roles. A program file is then parsed and analyzed against those specifications. Instead of running the program for its concrete result, the checker computes the abstract stack effect of each token and of each user-defined word. The same mechanism also produces source-aware diagnostics when adjacent effects cannot be composed. In that sense, the implementation follows the abstract-interpretation idea of executing programs over a finite abstract domain rather than over concrete values \cite{cousot1977abstract}.

Two design choices make the system concise and extensible. First, types are not hard-coded; they are loaded as a small finite lattice-like structure. Second, control flow is not built into the checker alone; it can be described declaratively with \texttt{syntax:} and \texttt{effect:} blocks. The mathematical content of the checker is therefore concentrated in a small effect algebra. This approach is informed by earlier work on algebraic specifications of stack effects, multiple stack effects, wildcard-based polymorphism, and compositional type systems for stack-oriented languages \cite{poial1991algebraic,poial1992multiple,stoddart1993type,poial2002wildcards,saabas2006compositional,poial2006typing}.

\section{Type Domain}
\subsection{Type hierarchy}
Let $\Types$ be a finite set of type names equipped with a precision order $\leq$ \cite{davey2002lattices}. We read $\tau_1 \leq \tau_2$ as ``$\tau_1$ is at least as precise as $\tau_2$,'' matching the sample declarations
\begin{lstlisting}
rel a-addr < c-addr
rel c-addr < addr
rel addr < x
rel flag < x
rel char < n
rel n < x
\end{lstlisting}
from the repository. Thus \texttt{a-addr} is more precise than \texttt{c-addr}, which is more precise than \texttt{addr}, and so on.

The implementation stores this relation as a finite transitive closure. Operationally, the checker only needs the meet of \emph{comparable} types.

\begin{definition}[Comparable meet]
For $\tau,\upsilon \in \Types$, define
\[
\tau \glb \upsilon =
\begin{cases}
\tau & \text{if } \tau \leq \upsilon,\\
\upsilon & \text{if } \upsilon \leq \tau,\\
\text{undefined} & \text{otherwise.}
\end{cases}
\]
\end{definition}

When the order is viewed as a finite lattice, this operation is the greatest lower bound on chains \cite{davey2002lattices}. Incomparable pairs have no usable meet for this checker and therefore cause a type clash.

\subsection{Stack symbols}
A \emph{stack symbol} is a type optionally annotated with a positive index, written $\tau$ or $\tau[i]$. The type denotes an abstract class of stack cells. The index denotes correlation across positions: if the same indexed symbol appears twice, those occurrences refer to the same abstract value constraint, not merely two values of the same coarse type.

For example, the bundled specification
\begin{lstlisting}
PLUS ( x[1] x[1] -- x[1] )
\end{lstlisting}
states that both arguments to \texttt{PLUS} and its result are tied to the same abstract type variable. Likewise,
\begin{lstlisting}
DUP ( x[1] -- x[1] x[1] )
\end{lstlisting}
states that duplication preserves the abstract identity of the top stack element.

Unnamed symbols are treated as implicitly fresh. After evaluation, the checker renumbers indices into a compact normal form, so effects are compared modulo alpha-renaming.

\section{Stack Effects}
\subsection{Syntax}
Let $\Sigma$ be the set of stack symbols. A stack effect is a pair of finite sequences
\[
\pi = (L \; -- \; R), \qquad L,R \in \Sigma^\ast.
\]
The sequence is written in ordinary Forth order: the rightmost element is the top of the stack. For example,
\[
(x[2]\;x[1]\; -- \; x[2]\;x[1]\;x[2])
\]
is the effect of \texttt{OVER}.

The empty effect
\[
\id = (\, -- \,)
\]
acts as the identity transformation.

\subsection{Normalization}
Two effects that differ only by a renaming of fresh indices should be considered equal. The implementation therefore normalizes inferred effects after every successful sequence evaluation. Symbols that occur only once and were never explicitly indexed are printed without an index; correlated or explicitly named symbols are assigned compact indices in a deterministic order.

This normalization step is not cosmetic. It lets the checker treat effects as abstract objects rather than artifacts of a particular evaluation order.

\section{Sequential Composition}
The primary operation of the checker is sequential composition of effects.

\begin{definition}[Boundary composition]
Let
\[
\pi_1=(L_1 -- R_1), \qquad \pi_2=(L_2 -- R_2).
\]
Their composition $\pi_1 \compose \pi_2$ is computed from the touching boundary between the output stack $R_1$ and the input stack $L_2$:
\begin{enumerate}
\item If $R_1$ is empty, the result is $(L_2L_1 -- R_2)$.
\item If $L_2$ is empty, the result is $(L_1 -- R_1R_2)$.
\item Otherwise let $a$ be the top element of $R_1$ and $b$ the top element of $L_2$.
  If $a.\mathsf{type} \glb b.\mathsf{type}$ is undefined, composition fails.
  If it is defined, replace both symbols everywhere by a fresh symbol whose type is that meet, remove the matched boundary pair, and continue recursively.
\end{enumerate}
\end{definition}

The fresh symbol keeps the stronger of the two type constraints. Since lower elements in the order are more precise, the meet is the more informative symbol. This is why a branch that expects \texttt{c-addr} can still compose with one that only proves \texttt{a-addr}: the merged requirement becomes \texttt{a-addr}.

The command-line example from the repository illustrates the rule:
\begin{lstlisting}
1 2 + DUP
\end{lstlisting}
with specifications
\begin{lstlisting}
literal integer ( -- n )
+ ( n n -- n )
DUP ( x[1] -- x[1] x[1] )
\end{lstlisting}
The checker prints
\begin{lstlisting}
1   ( --  n[2] )
2   ( --  n[2] )
+   ( n[2] n[2] --  n[1] )
DUP ( n[1] --  n[1] n[1] )
\end{lstlisting}
and concludes that the whole program has effect
\[
( -- n[1]\;n[1]).
\]

\begin{proposition}
Modulo normalization of fresh indices, compatible effects form a partial semigroup under $\compose$.
\end{proposition}

\noindent
This proposition explains why the implementation may evaluate a sequence left to right: whenever every adjacent composition is defined, the final result does not depend on accidental wildcard names. The terminology comes from semigroup theory, but the operation is partial because some boundary unifications fail \cite{howie1995semigroup}.

\section{Branching and Greatest Lower Bounds of Effects}
Sequential composition is not enough for control flow. A branch requires a single summary effect that safely represents multiple alternatives.

The checker therefore implements a greatest-lower-bound operation on whole effects. Intuitively, it merges branch summaries pointwise after accounting for stack cells that are preserved unchanged across both input and output. This is necessary because different branches may manipulate different depths of the stack while still agreeing on a common surrounding context.

More concretely, if one effect has $k$ more inputs than another, it must also have $k$ more outputs. The implementation then extracts a preserved prefix of length $k$ from the larger effect by requiring each corresponding input/output pair in that prefix to be comparable. After that prefix has been merged, the remaining active slices are unified pointwise. If any comparison fails, the branch summaries are incompatible.

This operator is used by the \texttt{either} constructor in declarative control semantics. The example
\begin{lstlisting}
: MYCHOICE IF ! ELSE C! FI ;
\end{lstlisting}
is inferred as
\begin{lstlisting}
MYCHOICE ( x a-addr flag -- )
\end{lstlisting}
because the two branches have effects
\[
(x\;a\text{-}addr -- )
\qquad\text{and}\qquad
(x\;c\text{-}addr -- ),
\]
and
\[
a\text{-}addr \glb c\text{-}addr = a\text{-}addr.
\]
The merged effect therefore preserves the stronger address requirement.

\section{Iteration and Idempotence}
Loops require a stable summary of repeated execution. The checker does not attempt a general fixed-point engine of the kind common in abstract interpretation \cite{cousot1977abstract}. Instead it uses a compact algebraic criterion based on idempotence.

\begin{definition}[Repeated effect]
For an effect $\pi$, define
\[
\pi^\star = \pi \glb (\pi \compose \pi),
\]
when both the composition and the greatest lower bound are defined.
\end{definition}

This is a one-step stabilization test. If one iteration and two iterations admit a common lower bound, that lower bound is taken as the loop summary. If no such lower bound exists, the body is rejected as a non-idempotent repeated effect.

This operator is used by the \texttt{repeat} constructor in control specifications. Algebraically, the checker is looking for an effect that is stable under another round of execution. The intended reading is not full Kleene closure, but rather a compact approximation of a loop invariant on the stack.

The bundled program
\begin{lstlisting}
: MYLOOP BEGIN SWAP OVER WHILE INVERT REPEAT ;
\end{lstlisting}
is inferred as
\begin{lstlisting}
MYLOOP ( flag[1] flag[1] -- flag[1] flag[1] )
\end{lstlisting}
which is precisely an idempotent invariant: repeating the body preserves the same abstract pair of flags.

\section{Declarative Specifications}
\subsection{Ordinary words, literals, and parser words}
The \texttt{specs} file associates stack effects with words and adds parsing metadata. Representative examples are:
\begin{lstlisting}
OVER ( x[2] x[1] -- x[2] x[1] x[2] )
@ ( a-addr -- x )
CONSTANT parse word define ( x -- )
VARIABLE parse word define ( -- a-addr )
literal integer ( -- n )
\end{lstlisting}

Additional clauses refine the operational role of a word.
\begin{itemize}
\item \texttt{parse word} means the word consumes the next token from the source.
\item \texttt{parse until DELIM} means the word consumes source text until a delimiter.
\item \texttt{define}, \texttt{define colon}, \texttt{define constant}, and \texttt{define variable} mark defining words.
\item \texttt{control ROLE} assigns a control-flow role such as \texttt{if}, \texttt{else}, \texttt{fi}, \texttt{begin}, or \texttt{while}.
\item \texttt{immediate}, \texttt{state interpret}, and \texttt{state compile} constrain where the word is legal.
\end{itemize}

Thus the checker is not only typing tokens; it also models the source-level behavior needed to parse Forth-like syntax in a way that aligns with standard Forth parsing disciplines \cite{forth2012standard}.

\subsection{Control structures}
Control structures are specified with paired \texttt{syntax:} and \texttt{effect:} blocks. This gives a compact declarative account of source-level control words, a topic that has also been studied in formal models of Forth control structures \cite{power2004formal}. For example, the sample repository contains:
\begin{lstlisting}
syntax: IF <then branch> [ELSE <else branch>] FI
  effect:
    IF
    either <then branch> <else branch>
\end{lstlisting}

and
\begin{lstlisting}
syntax: BEGIN <loop prefix> WHILE <loop body> REPEAT
  effect:
    repeat <loop prefix> WHILE
    repeat <loop body>
\end{lstlisting}

The \texttt{syntax:} clause describes a control-language pattern using concrete boundary words and metavariables for internal segments. Optional parts are written with square brackets. The \texttt{effect:} block is then compiled into an expression language with the grammar
\[
E ::= \eps \mid \langle s \rangle \mid c \mid E\,E \mid \textsf{either}(E,E) \mid \textsf{repeat}(E),
\]
where $\langle s \rangle$ denotes a named program segment and $c$ denotes the runtime effect of a control word such as \texttt{IF}, \texttt{WHILE}, or \texttt{DO}.

The denotational semantics is:
\begin{align*}
\sem{\eps} &= \id,\\
\sem{\langle s \rangle} &= \text{the inferred effect of segment } s,\\
\sem{c} &= \text{the declared runtime effect of control role } c,\\
\sem{E_1 E_2} &= \sem{E_1} \compose \sem{E_2},\\
\sem{\textsf{either}(E_1,E_2)} &= \sem{E_1} \glb \sem{E_2},\\
\sem{\textsf{repeat}(E)} &= \sem{E}^\star.
\end{align*}

This design cleanly separates parsing from typing. Once a source fragment has been segmented according to a \texttt{syntax:} declaration, its meaning is calculated entirely by the effect algebra above.

\section{Inference of User-Defined Words}
At top level, the checker interprets defining words immediately. A colon definition starts a new sequence, parses until its terminating control word, computes the effect of the body, and installs the resulting effect as the specification of the new word. Constants and variables are handled similarly but with special defining shapes:
\[
\texttt{CONSTANT} : (x --), \qquad \texttt{VARIABLE} : (-- a\text{-}addr).
\]

The implementation therefore supports the standard Forth workflow in which the specification dictionary grows during analysis \cite{forth2012standard}. Inferred definitions are written to a log file; for the sample program the checker records:
\begin{lstlisting}
MYLOOP   ( flag[1] flag[1] -- flag[1] flag[1] )
MYCHOICE ( x a-addr flag -- )
MYTEST   ( a-addr[2] x[1] flag -- a-addr[2] x[1] a-addr[2] )
MYFI     ( x[1] x[1] flag -- x[1] x[1] )
MYAGAIN  ( a-addr[1] -- a-addr[1] )
MYSEQ    ( a-addr[2] x[1] -- x[1] a-addr[2] a-addr[2] )
MYSEQ2   ( a-addr[1] a-addr[1] -- )
\end{lstlisting}

These examples demonstrate several features of the checker at once.
\begin{itemize}
\item \texttt{MYCHOICE} uses branch merging over comparable address types.
\item \texttt{MYTEST} shows prefix-aware branch merging when the alternatives manipulate different active depths of the stack.
\item \texttt{MYLOOP} and \texttt{MYAGAIN} show loop summaries computed through repeated-effect stabilization.
\item \texttt{MYSEQ2} uses the polymorphic example word \texttt{PLUS}, showing that indexed symbols propagate through a larger composition.
\end{itemize}

\section{Diagnostics and Practical Scope}
The checker is designed as a static evaluator rather than a full runtime. It reports source spans for malformed specifications, unknown types, missing parser delimiters, illegal control words, and type clashes in both top-level programs and definitions. When composition fails, the checker reports the surrounding context and the conflicting symbols at the composition boundary.

The current implementation intentionally stays small. It relies on the declared specification set, rejects nested defining words inside definitions, enforces interpretation-state restrictions, and treats the external type hierarchy as the semantic authority. These restrictions are acceptable for a checker whose purpose is to validate stack discipline and infer word summaries, not to execute arbitrary Forth code.

\section{Conclusion}
The type checker described in this paper shows that a useful static analysis for a Forth(-like) language can be built from a very small mathematical core. A finite type lattice supplies precision information; indexed stack symbols express correlations; sequential composition forms a partial semigroup of effects; branch summaries are greatest lower bounds of effects; and loop summaries are accepted when repeated execution stabilizes to an idempotent effect. Because the checker loads both the type hierarchy and the word specifications from external files, the same implementation can be retargeted to other vocabularies and dialects with little or no change to the core.

In that sense, the paper's main lesson is methodological: many of the dynamic-looking features of a concatenative language become statically tractable once they are phrased as operations in an effect algebra.

The type checker successfully analyzed its own source using the Forth 2012 standard type system \cite{forth2012standard}.

\bibliographystyle{plain}
\bibliography{typechecker-paper}

\end{document}
